% =====================================================================
% When Errors Become Narratives: A Longitudinal Taxonomy of Silent
% Failures in a Production LLM Agent Runtime
%
% IEEE Software resubmission (SW-2026-06-0312) — compliant experience
% report version (NOT the full arXiv draft, which exceeded the 4,200-word
% limit and triggered the desk-return).
%
% IEEE journal format (IEEEtran). Build: pdflatex ieee_software_manuscript.tex
%   (run twice for cross-references). No bibtex: bibliography is inline.
% No external files, no figures — 2 tables only. Standard TeXLive packages.
%
% Word budget (IEEE: <=4,200 incl. 250/table+figure; 2 tables = 500):
%   body prose ~3,100 + 500 = ~3,600 counted words (<4,200). Abstract 148.
%
% Author: Wei Wu <wuweinanonuaa@gmail.com>, Independent Researcher.
% ORCID: 0009-0009-1176-7817.
% =====================================================================
\documentclass[10pt,journal]{IEEEtran}
\IEEEoverridecommandlockouts

\usepackage[T1]{fontenc}
\usepackage[utf8]{inputenc}
\usepackage{cite}
\usepackage{amsmath,amssymb}
\usepackage{booktabs}
\usepackage{array}
\usepackage{xcolor}
\usepackage[hidelinks]{hyperref}

\newcommand{\failplausible}{\emph{fail-plausible}}
\newcommand{\sysname}{\texttt{openclaw-model-bridge}}

\begin{document}

\title{When Errors Become Narratives: A Longitudinal Taxonomy of Silent
Failures in a Production LLM Agent Runtime}

\author{Wei~Wu%
\thanks{W.~Wu is an independent researcher.
E-mail: wuweinanonuaa@gmail.com; ORCID: 0009-0009-1176-7817.
A companion extended preprint is available as arXiv:2606.14589.}}

\markboth{IEEE Software --- Feature Article (Experience Report)}%
{Wu: When Errors Become Narratives}

\maketitle

\begin{abstract}
Large language model (LLM) agent runtimes increasingly run unattended for
weeks, calling tools and pushing results to humans. I report an eight-week
longitudinal study of \emph{silent failures} in one such production system:
a personal-assistant agent runtime defended by 4{,}286 unit tests and 827
governance checks, in which I documented 22 incidents with full root-cause
postmortems. One failure class proved unique to systems that \emph{speak}:
the LLM transforms an upstream error into fluent, plausible, false output
for the user---a behavior I call \textbf{fail-plausible}.
Three findings unsettle common assumptions: $\sim$70\% of silent failures
were caught by a human reading the product, not the automated stack; the
governance audit prevented 0\% of novel incidents but blocked 87\% from
recurring; and the longest failures lived in the \emph{seams} between simple
parts, not inside complex ones. I distill the defenses that worked into
practices adoptable without my stack.
\end{abstract}

\begin{IEEEkeywords}
LLM agent systems, software reliability, silent failures, fail-plausible,
observability, fault taxonomy, postmortem analysis, hallucination,
AI-assisted software engineering, experience report.
\end{IEEEkeywords}

\IEEEpeerreviewmaketitle

\section*{Three Practitioner Takeaways}
\begin{enumerate}
\item \textbf{Schedule user-view observation as a first-class observability
signal.} In this study a human reading the actual pushed output caught
$\sim$70\% of silent failures that thousands of green tests and hundreds of
green checks missed. Put 30 minutes a week on the calendar to \emph{use} the
system, not test it.
\item \textbf{Keep diagnostics out of data channels, and preserve error cause
across every hop.} An LLM downstream will narrate whatever lands in its
context. A single \texttt{>\&2} redirection severed one fabricated ``platform
crisis.'' Treat stderr discipline and cause-preserving error chains as
security controls for agent pipelines.
\item \textbf{Convert each incident into a mechanized check, and prove every
guard by sabotage.} A lesson that stops at a point fix recurred within two
days; an unvalidated guard is indistinguishable from a vacuous one (I found
67 checks silently running empty strings for months).
\end{enumerate}

\section{A Crisis That Never Happened}
One morning my agent runtime pushed me a confident industry analysis about a
``Hugging Face platform crisis.'' It was well written, on schedule,
thematically plausible---and completely fabricated. No alarm fired. Every unit
test was green. Every health check passed.

Here is what actually happened. A nightly knowledge-synthesis job scraped
content containing an isolated Unicode surrogate, which made
\texttt{json.dump} raise mid-write. The truncated request body drew an HTTP
400 from the model adapter. So far, an ordinary bug. The damage came next: the
job's logging function wrote its diagnostics to \emph{stdout}, and the caller
captured stdout by shell command substitution as the \emph{signal payload}.
The cache filled with strings like \texttt{HTTP Error 400: Bad JSON \ldots
waiting 3s before retry}. The downstream reduce-step LLM, prompted to find
cross-domain signals, did exactly what language models do with
anomalous-but-thematic text: it composed a fluent narrative---platform trouble
being the most probable story for error-code vocabulary---and the system
delivered it to me as a routine insight digest. Every component succeeded. The
pipeline laundered an encoding bug into industry analysis.

The fix that severed the chain was a single character: \texttt{>\&2},
redirecting diagnostics to stderr so they could never enter the data channel.
I added defense in depth around it---surrogate sanitization, an encoding-error
policy, anti-fabrication prompt guards---but the load-bearing repair was one
redirection operator.

This incident is not an outlier. It is the most vivid member of a failure
class I believe is specific to LLM agent systems, and it changed how I think
about reliability for software that generates language.

\section{Fail-Plausible: The Silent Failure of Systems That Speak}
The reliability literature has long known that the most damaging production
failures are not crashes. \emph{Gray failure} degrades cloud systems while
their failure detectors report health~\cite{ref1}; \emph{fail-slow} hardware
throttles for hours before anyone suspects the disk~\cite{ref2}. Both share
one defining property---\textbf{differential observability}: the application
suffers, but the observer built to notice does not.

LLM agent runtimes inherit this entire problem class and add a worse one. An
agent runtime is a generator of fluent language. When an upstream error leaks
into its context window, its failure mode is not silence; it is \emph{plausible
speech}. I call this \textbf{fail-plausible}: a failure in which the system
transforms an internal error into coherent, contextually appropriate, and
false output. Fail-plausible is to agent systems what gray failure was to
cloud infrastructure---the dominant, hardest-to-see class---but it is strictly
worse for the observer. Gray failure starves the detector of a signal;
fail-plausible hands the human a counterfeit one. The observer is not merely
blind; it is being convincingly lied to by the failure itself.

Crucially, the model is not malfunctioning. In every fabrication I documented,
the LLM behaved exactly as trained: fluent completion over its context. The
failure is a \emph{systems} failure---the pipeline delivered polluted context
(captured error logs, stale alerts, unlabeled enrichment) to a component whose
job is to narrate whatever it is given. That reframing matters because it moves
the defense off the model and onto the system, where engineers actually have
leverage.

\section{The System, and How I Studied It}
The subject is a two-layer middleware (public repository \sysname) connecting
self-hosted and commercial LLMs to an open-source personal-agent framework on
WhatsApp and Discord. It has run continuously on a single macOS host since
March 2026, with roughly 40 scheduled jobs across two cron schedulers, 8 LLM
providers, a tool-governance proxy, and a knowledge-base memory
plane---defended at the study cutoff by 4{,}286 unit tests in 121 suites and
827 declarative governance checks. One human operator (me) and one AI
engineering collaborator (Claude, used through a coding-agent interface)
develop and operate it; the runtime itself is powered by separate Qwen3-class
models.

The corpus is every incident between 2026-04-09 and 2026-06-02 that reached
production, had a \emph{silent phase} (failing while all automated indicators
stayed green), and was closed with a full postmortem. Twenty-two qualify. Each
follows a mandatory in-repo protocol that requires, \emph{before any fix}: a
full causal-chain diagram (timeline $\times$ layer $\times$ logic); a
three-layer root cause (trigger / amplifier / concealer); a timeline
reconstruction; a ``why now, not before'' condition-combination analysis; and
an entry feeding the governance ontology. Every postmortem is a standalone
public document, and the numbers in this article are mechanically derivable
from the repository---which I offer as the primary check on narrative bias,
including the bias of an AI-assisted author writing about AI-assisted
operation.

I classify by \textbf{mechanism} (how the error evaded observation), not
\textbf{location} (which job failed), because the same mechanism recurs across
unrelated components: a location taxonomy had no defensive value, while one
mechanism-level scanner immunized every location at once.

\section{Five Ways an Agent Fails in Silence}

\begin{table*}[t]
\centering
\caption{A mechanism-oriented taxonomy over 22 incidents.}
\label{tab:taxonomy}
\footnotesize
\begin{tabular}{@{}p{2.6cm}p{6.6cm}p{1.7cm}p{4.4cm}@{}}
\toprule
\textbf{Class} & \textbf{Mechanism} & \textbf{Silence span} & \textbf{Signature} \\
\midrule
A --- Environment/platform quirk & Logic is correct; the runtime
environment's implicit behavior defeats it & hours--weeks & green in dev,
silent in prod \\
\addlinespace
B --- Design-assumption mismatch & Code assumes a deployment topology,
contract, or input shape that reality violates & days & tests mirror the
assumption, not the caller \\
\addlinespace
C --- Error swallowing \& dilution & The error is eaten by a layer, or
stripped of cause across layers & hours--days & the alert that arrives carries
no usable information \\
\addlinespace
D --- Chained hallucination (\textbf{fail-plausible}) & The LLM converts
polluted context into confident false output & hours--days & the user receives
counterfeit health \\
\addlinespace
E --- Operational omission \& forensic blind spot & A deploy/registration step
never happened, or the diagnostic instrument itself is blocked and reads as
``normal'' & days--\textbf{60 days} & declared state $\neq$ runtime state;
instruments lie \\
\bottomrule
\end{tabular}
\end{table*}

Three classes will be familiar to systems engineers (Table~\ref{tab:taxonomy}).
\textbf{Class A} is the dev--prod environment gap with teeth: bash 3.2 not
propagating an \texttt{ERR} trap into a function silently disarmed a watchdog;
BSD \texttt{awk} aborting on invalid UTF-8 killed the \emph{monitoring} script
for seven days. \textbf{Class C} is error dilution: an evening digest failed
with ``HTTP 502'' for two days while the true cause---a fallback provider's
quota exhausted after the primary's circuit breaker opened---sat in an upstream
response body that three successive layers each declined to read.
\textbf{Class E} is the largest and longest: a new job fully implemented,
tested, and registered simply never ran because the final step (writing the
crontab line) was a human memory item, and three small bugs conspired to
report success.

Two classes are sharper in an agent runtime. \textbf{Class B} includes any
\emph{positional} parse of LLM output: a parser that indexed
\texttt{lines[i+1]}, \texttt{lines[i+2]} shifted every field by one the day the
model omitted a line---because instruction-following is a distribution, not a
contract. And \textbf{Class D} is fail-plausible, narrated above and
elaborated below.

The fabrications were not all dramatic. A weekly review job whose LLM call
failed fell back to a mechanical line filter that emitted leftover container
headings as ``review content'' and wrote \texttt{"llm": true} to its status
file unconditionally---\emph{a hallucination implemented in shell.} An evening
digest, handed the day's papers as unlabeled context, inferred my project
``must have shipped'' and announced a release of an internal version number
that exists only in a changelog. The common structure is a \textbf{pollution
chain}: a Class A/B/C failure deposits non-signal content where a downstream
LLM expects signal; the LLM completes fluently; and the output inherits the
\emph{form} of health with the \emph{content} of failure.

\section{Three Uncomfortable Findings}
\textbf{A human reading the product out-detected the entire automated stack.}
Sorting incidents by discovery channel, roughly 70\% were ultimately caught by
\emph{user-view observation}---me noticing ``this digest looks shallow,'' ``why
two windows?'', ``I didn't get yesterday's analysis.'' The 4{,}286 tests, 827
checks, and a 19-point preflight stayed green through most incidents; unit
tests caught essentially none of this corpus (by selection---anything they
caught never became silent). I institutionalized a weekly 30-minute observation
ritual: no coding, just using the system as a user along four axes (alert
noise, push latency, information density, response quality). It kept
out-detecting the automated stack. The research-grade version of this---an LLM
``observer'' that critiques yesterday's output~\cite{ref9}---is promising and
found real regressions, but it is itself an LLM component subject to every
class in this taxonomy (mine shipped with a Class B bug and a hallucinated
truncation), so it needs the same governance as what it judges.

\textbf{The audit prevented nothing and blocked almost everything.} Mid-study I
audited my own defense framework against the first 15 incidents with three
questions each: could the audit, as it existed \emph{before} the incident, have
caught it (Q1)? Why not (Q2)? Do the guards added \emph{after} block the class
going forward (Q3)?

\begin{table}[t]
\centering
\caption{Auditing the defense framework against the first 15 incidents.}
\label{tab:audit}
\begin{tabular}{@{}p{6.0cm}r@{}}
\toprule
\textbf{Metric} & \textbf{Value} \\
\midrule
Ex-ante prevention (Q1) & \textbf{0 / 15 = 0\%} \\
Ex-post regression blocking (Q3) & \textbf{13 / 15 = 87\%} \\
Misses rooted in a dimension never conceived & 12 / 15 = 80\% \\
\bottomrule
\end{tabular}
\end{table}

The 0\% is not an indictment---it is the job description. Audits, like
regression suites, encode the past; 80\% of misses were \emph{blank categories}
no invariant had ever contemplated. The honest posture: maximize the regression
rate, accept that prevention of \emph{novel} classes comes from elsewhere
(user-view observation, adversarial review, target-environment exposure), and
minimize the \emph{conversion latency} from a new incident to a mechanized
guard.

\textbf{The longest failures lived in seams, not in complex code.} Latency
tracked the failure mechanism, not code complexity. Code-level bugs die
young---tests and the next run catch them. What survived for weeks lived where
no test runs: deployment topology, OS policy, the gap between declared and
runtime state, and monitoring-of-monitoring. The corpus's longest silence (60
days) was a backup failing with \texttt{EPERM}, where every forensic tool I
reached for was \emph{itself} being silently denied by the macOS TCC sandbox
and returning empty output the pipeline recorded as ``normal.'' An instrument
that cannot tell ``nothing there'' from ``I was not allowed to look''
manufactures false reassurance. Latency, in other words, measures
\emph{observational distance}: the further a mechanism sits from any existing
observer, the longer it lives. I now treat silence-latency percentiles as a
reliability metric in their own right---for silent failures,
time-to-\emph{detect} dominated time-to-repair by one to two orders of
magnitude.

\section{What Actually Defends Against Silent Failure}
The defenses that held up are less a framework than a discipline. Five
practices, mapped to the takeaways above.

\textbf{Make diagnostics structurally unable to become data, and never let an
error lose its cause.} Shell diagnostics go to stderr, enforced
repository-wide by a scanner; alerts are tagged at every producer and stripped
from chat context \emph{before} truncation; error chains must preserve upstream
cause across every hop. In an agent system the ``client'' of an error is often
another LLM prompt, so a diluted or misrouted error is one step from becoming
fabrication input. This is the Class D defense, and it lives mostly in
plumbing, not prompts.

\textbf{Treat the LLM's context as a supply chain.} Beyond hygiene, label
provenance: a source-credibility module tags every ingested source on a
five-tier scale at prompt-injection time, and a shared anti-fabrication guard
module (imported, never copy-pasted) injects a cumulative ladder of
prohibitions into all LLM-calling jobs---the upper levels containing
\emph{literal} phrases harvested from real fabrication incidents, including the
exact false release string from the digest above. After deployment, the
targeted fabrication patterns (multi-hop causal chains, ``therefore''-style
necessity claims) dropped 53--92\% in the affected job. Prompt guards are the
\emph{last} layer, behind hygiene and provenance---never the first.

\textbf{Climb the three-step ladder: point fix $\rightarrow$ meta-rule
$\rightarrow$ scanner.} A point fix repairs this bug; empirically
insufficient---an import-omission fix at one site recurred at eight sites via an
automated injector two days later, because the lesson existed only as a diff. A
meta-rule names the lesson as a cross-case rule (I reached 23:
diagnostics-to-stderr, key-based parsing, alert-path independence,
declared-state convergence, reserved-file unwritability). A mechanized scanner
encodes the meta-rule as a check that runs in CI and the daily audit (14 by
study end). Only at step three does recurrence become structurally impossible
rather than culturally discouraged. Every meta-rule that reached step three has
zero recorded recurrences; every recurrence I logged involved a lesson that
stopped at step one.

\textbf{Prove every guard by sabotage.} New guards must be shown \emph{alive}---
introduce the violation they target, watch them fire, revert. This caught
guards that matched their own assertion strings, tests whose fixtures mirrored
the wrong assumption, and 67 governance checks that had been executing empty
strings for months (discovered only when a \emph{new} invariant's sabotage
refused to fail). In a regime whose primary failure is silence, an unvalidated
detector is indistinguishable from a vacuous one.

\textbf{Make alarms outlive their subject, and make monitoring watch itself.} A
nine-hour total outage went unannounced because three independent alarms each
failed---including a quiet-hours filter that suppressed the very channel meant
for emergencies. The rule: an alert path must not depend on the failing subject
(gateway-down notices travel a transport the gateway cannot take down), and the
watcher needs a watcher (an independent job age-checks a heartbeat canary
file).

\section{Seams, Not Components}
Facing a week with five incidents, I asked the obvious question: \emph{is this
system failing because it has become too complex?} The postmortem-backed answer
was more useful: \textbf{no individual failing part was complex.} A symlink. An
\texttt{abspath} call. A one-line registry entry. A boolean default. Every part
was simple and locally correct; every incident lived in a \emph{combination}.
Components are covered by thousands of tests; \textbf{combinations grow
superlinearly, and tests cover only the combinations someone imagined.}
Complexity is about parts; incidents are about seams.

This inverts the reflex response to an incident---\emph{add a guard}---because a
guard is itself a part, and therefore new seams. (A convergence engine I built
to close the declared-vs-runtime seam opened an observer-mutates-observed seam
that caused three incidents before I caged it with a rule: \emph{audit observes
and never mutates}.) The posture I landed on, codified as a standing ``Sunset
Law'': before adding any mechanism, try to \emph{retire} an equivalent one; one
logical entity should have one physical representation (multiple representations
\emph{will} drift, and the bug lives between them); and defenses are themselves
incident surfaces---prefer \textbf{seam reduction} (unify representations,
shrink the dev--prod gap, keep observers read-only) over defense accretion. In
the weeks since, the system retired more representation duplicates than it added
invariants, and---with due caution about confounds---incident frequency fell
while feature velocity did not.

One honest note on method. A single human producing 22 publication-grade
postmortems, 90 invariants, and 14 scanners in eight weeks is achievable
\emph{only} with an AI collaborator---and that same collaborator will, unguided,
pattern-match symptoms to plausible fixes at high speed, industrializing the
cosmetic trigger-level fix (one documented chain was five cascading ``fixes'' to
a problem whose correct resolution was a single file copy). The rigid postmortem
protocol exists largely to counter that. AI collaboration shifts the binding
constraint of reliability engineering from implementation bandwidth to
\emph{judgment discipline}---which is exactly what the meta-rules encode. This
is a single-system case study, and I claim the \emph{classes and structures}
generalize to long-running agent runtimes with human-facing output, not the
frequencies; the public repository and a number-by-number data inventory are
offered so others can check the chain and, ideally, replicate the method on
their own systems.

\section{The Failure That Should Frighten You}
Eight weeks of complete postmortems yield a five-class taxonomy of silent
failures, of which one---fail-plausible chained fabrication---is specific to
systems that speak. The record is uncomfortable in useful ways: thousands of
tests and hundreds of checks prevented none of the novel incidents while
blocking 87\% from recurring; the best detector was a human reading the product;
and the longest failures lived in the seams between simple, correct parts. The
failure your agent system should frighten you with is not the crash. It is the
confident paragraph, on schedule, in perfect grammar, about a crisis that does
not exist---pushed to a user who has every reason to believe it, by a pipeline
in which every component worked.

\section*{Acknowledgment of AI Assistance}
This article and the underlying postmortems were prepared with the assistance
of Anthropic's Claude, used as an engineering and writing collaborator through
a coding-agent interface. Claude contributed to incident triage, postmortem
drafting, repository data extraction, and prose editing. The author, Wei Wu, is
the sole author, defined all research questions and conclusions, verified every
reported figure against the public repository, and is solely responsible for the
content. No generative-AI system is a listed author. (See the project's
disclosure statement for the full division of contributions.)

\begin{thebibliography}{12}
\bibitem{ref1} P. Huang et al., ``Gray Failure: The Achilles' Heel of
Cloud-Scale Systems,'' \emph{HotOS}, 2017.
\bibitem{ref2} H. S. Gunawi et al., ``Fail-Slow at Scale: Evidence of Hardware
Performance Faults in Large Production Systems,'' \emph{USENIX FAST}, 2018.
\bibitem{ref3} M. Cemri et al., ``Why Do Multi-Agent LLM Systems Fail?'' (MAST),
arXiv:2503.13657, 2025.
\bibitem{ref4} B. Ranganathan, M. Zhang, and K. Wu, ``Enhancing Reliability in
AI Inference Services: An Empirical Study on Real Production Incidents,''
arXiv:2511.07424, 2025.
\bibitem{ref5} C. Ezell, X. Roberts-Gaal, and A. Chan, ``Incident Analysis for
AI Agents,'' \emph{AAAI/ACM AIES}, 2025; arXiv:2508.14231.
\bibitem{ref6} B. Beyer, C. Jones, J. Petoff, and N. R. Murphy (eds.),
\emph{Site Reliability Engineering: How Google Runs Production Systems},
O'Reilly, 2016.
\bibitem{ref7} A. Basiri et al., ``Chaos Engineering,'' \emph{IEEE Software},
vol. 33, no. 3, pp. 35--41, 2016.
\bibitem{ref8} S. Ghosh, M. Shetty, C. Bansal, and S. Nath, ``How to Fight
Production Incidents? An Empirical Study on a Large-Scale Cloud Service,''
\emph{ACM SoCC}, 2022.
\bibitem{ref9} L. Zheng et al., ``Judging LLM-as-a-Judge with MT-Bench and
Chatbot Arena,'' \emph{NeurIPS}, 2023; arXiv:2306.05685.
\bibitem{ref10} L. Huang et al., ``A Survey on Hallucination in Large Language
Models: Principles, Taxonomy, Challenges, and Open Questions,''
arXiv:2311.05232, 2023.
\bibitem{ref11} P. Runeson and M. H\"ost, ``Guidelines for Conducting and
Reporting Case Study Research in Software Engineering,'' \emph{Empirical
Software Engineering}, vol. 14, no. 2, pp. 131--164, 2009.
\bibitem{ref12} M. Taraghi, M. M. Morovati, and F. Khomh, ``Real Faults in
Model Context Protocol (MCP) Software: A Comprehensive Taxonomy,''
arXiv:2603.05637, 2026.
\end{thebibliography}

\begin{IEEEbiographynophoto}{Wei Wu}
is an independent researcher focused on the reliability of large language model
agent runtimes---silent-failure taxonomy, production observability, and the
governance of tool-calling agents. He designs, operates, and studies
\texttt{openclaw-model-bridge}, the production personal-assistant agent runtime
from which the incidents in this article are drawn. Contact him at
wuweinanonuaa@gmail.com; ORCID 0009-0009-1176-7817.
\end{IEEEbiographynophoto}

\end{document}
